\renewcommand{\tSoundText}{MSX PSG for music and effects.}
\renewcommand{\tVideoText}{A mandatory \textbf{V9990 cartridge} for all screens after the initial check, including credits, menu and gameplay.}
\renewcommand{\tMonitorText}{Access to both MSX and V9990 video outputs: two displays or a monitor with input selection.}
\renewcommand{\tControlsText}{Cursor keys; Space or M to jump. Joystick in either port, with either button.}
\renewcommand{\tSetMonitor}{Connect both video outputs. Begin with the MSX output to check detection, then switch to V9990 for the rest of the game.}
\renewcommand{\tSetOn}{Switch the MSX on. Its output shows \texttt{V9990 found!}. If it says \texttt{V9990 NOT found!}, switch off and check the V9990 cartridge.}
\renewcommand{\tBlackScreen}{\textbf{If you cannot see the picture}, check the selected input against section 3. Also check the cartridges and RAM. The supplied launcher switches video outputs automatically in the emulator.}
\renewcommand{\tCreditsTitle}{9. Credits and rights}
\renewcommand{\tSourceNote}{\href{https://gitlab.com/dihalt/madmix2-v9990}{Source code available on GitLab.} \textbf{Please improve the game if you feel so!}}
\newcommand{\tCreditsExtra}{\item DiHalt Studio opening logo and fanfare: DiHalt Studio presentation assets.\item Opening illustration: artwork supplied for this edition. Title logo, characters and game music: original game material.}
\renewcommand{\tSetJoy}{Connect a joystick to port 1 or 2 if you want to use one.}
\renewcommand{\tSetGame}{Insert the game cartridge in the other slot. For MegaFlashROM SD SCC+, copy the ROM to the SD as \texttt{MADMIX2.ROM} and flash with \texttt{OPFXSD MADMIX2.ROM /A16}. \texttt{/U} is optional: without it, the game uses a subslot; with it, it uses the primary slot. \texttt{/A16} forces ASCII16 if automatic detection fails.}
